\scanpage{73}
\margin{\unsure{$d$|$J$|$\delta$}}

\[
\boxed{\mathrm{sn}^{\prime\,2}z=(1-z^2)(1-k^2z^2)}\qquad \mathrm{sn}\,0=0\quad \mathrm{sn}'0=1 \tag*{(1)}
\]
\[
\boxed{\mathrm{sn}^2z=\frac{1}{\unsure{\wp|p}z-e_1}=\frac1{k^2}\bigl[\wp(z-\omega_1)-e_1\bigr]}\tag*{(2)}
\]
\[
\left\{\begin{aligned}
e_1&=-\frac{1+k^2}{3}\\
e_2&=\frac{2-k^2}{3}\\
e_3&=\frac{2k^2-1}{3}
\end{aligned}\right.
\qquad\begin{aligned}e_2-e_1&=1\\e_3-e_1&=k^2\end{aligned}
\]
\[
\boxed{\mathrm{sn}\,z=\sigma\omega_1\,e^{\eta_1 z}\frac{\sigma z}{\sigma(z+\omega_1)}
=-\sigma\omega_1\,e^{-\eta_1z}\frac{\sigma z}{\sigma(z-\omega_1)}}\tag*{(\unsure{1|3})}
\]
Si on choisit \uline{$k=-\frac{e^{\eta_1\omega_1}}{\sigma^2\omega_1}$} alors
\[
\boxed{\mathrm{sn}\,z\,\mathrm{sn}(z+\omega_1)=\frac1k}\tag*{(4)}
\]
\[
\left\{\begin{aligned}
\mathrm{sn}(z+2\omega_1)&=\mathrm{sn}\,z\\
\mathrm{sn}(z+2\omega_2)&=\mathrm{sn}\,z\\
\mathrm{sn}(z+2\omega_3)&=-\mathrm{sn}\,z\\
\mathrm{sn}(z+\deleted{2}4\omega_2)&=\mathrm{sn}\,z\\
\mathrm{sn}(z+4\omega_3)&=\mathrm{sn}\,z\\
\mathrm{sn}\,{-z}&=-\mathrm{sn}\,z
\end{aligned}\right.\tag*{(5)}
\]
$2\omega_1,4\omega_2$ et $4\omega_3$ sont périodes

$\mathrm{sn}\,z$ est impaire

La substitution $z\to z+2\omega_2$ ou $z+2\omega_{\unsure{3|2}}$ donne un changement de signe

\[
\left\{\begin{aligned}
&\text{\uuline{poles} : les homologues de }\text{\uuline{$\omega_1$}},\ \text{résidu }\frac1k\ ;\ \text{et }\omega_1+2\omega_2,\ \text{résidu }-\frac1k\\
&\text{\uuline{zéros} : les homologues de }\text{\uuline{$0$}},\ \text{dérivée }+1\ ;\ \text{et }2\omega_2,\ \text{dérivée }-1
\end{aligned}\right.\tag*{6}
\]
en choisissant convenablement $\omega_{2}$ :
\[
\underline{\mathrm{sn}\,\omega_2=1\quad\mathrm{sn}'\omega_2=0}
\qquad\underline{\mathrm{sn}\,\omega_3=\frac{-1}{k}\quad\mathrm{sn}'\omega_3=0}
\]

% archive sheet number: 72
